\documentclass[11pt]{article}
\usepackage[a4paper,margin=1in]{geometry}
\usepackage{booktabs}
\usepackage{tabularx}
\usepackage{graphicx}
\usepackage{float}
\usepackage[section]{placeins}
\usepackage{array}
\usepackage{amsmath}
\usepackage{xcolor}
\usepackage[authoryear,round,longnamesfirst]{natbib}
\bibpunct{(}{)}{;}{a}{}{,}
\usepackage[hidelinks]{hyperref}
\usepackage{fancyhdr}
\usepackage[T1]{fontenc}
\usepackage{lmodern}

\definecolor{faissblue}{HTML}{1F4D78}
\pagestyle{fancy}
\fancyhf{}
\lhead{\textit{faissR} supplementary material}
\rhead{\thepage}
\setlength{\headheight}{14pt}
\setlength{\emergencystretch}{3em}
\renewcommand{\headrulewidth}{0.4pt}

\newcommand{\code}[1]{\texttt{#1}}
\newcommand{\pkg}[1]{\textsf{#1}}
\newcommand{\proglang}[1]{\textsf{#1}}
\newcommand{\class}[1]{\texttt{#1}}
\newcolumntype{Y}{>{\raggedright\arraybackslash}X}
\renewcommand{\thetable}{S\arabic{table}}
\renewcommand{\thefigure}{S\arabic{figure}}

\title{\textbf{Supplementary Material}\\[0.4em]
\large \pkg{faissR}: Nearest-Neighbor Search with FAISS and cuVS in R}
\author{Moussa Kassim Mohamed$^{1,2,\dagger}$, Martin Ocharo$^{1,2,\dagger}$\\
Dalia Ahmed$^{1}$, Dupe Ojo$^{1}$\\
Alessia Vignoli$^{3,4}$, Leonardo Tenori$^{3,4}$\\
Dinesh Gupta$^{5}$, Silvano Piazza$^{6,7}$\\
Stefano Cacciatore$^{1,2,\ast}$}
\date{15 September 2026}

\begin{document}
\raggedbottom
\maketitle

\noindent
$^{1}$ Bioinformatics Unit, International Centre for Genetic Engineering and
Biotechnology (ICGEB), Cape Town 7925, South Africa.\\
$^{2}$ Department of Integrative Biomedical Sciences, Institute of Infectious
Disease \& Molecular Medicine (IDM), University of Cape Town, Cape Town 7925,
South Africa.\\
$^{3}$ Department of Chemistry ``Ugo Schiff'', University of Florence, Sesto
Fiorentino, Italy.\\
$^{4}$ Magnetic Resonance Center (CERM), University of Florence, Sesto
Fiorentino, Italy.\\
$^{5}$ Translational Bioinformatics Group, International Centre for Genetic
Engineering and Biotechnology (ICGEB), New Delhi, India.\\
$^{6}$ Computational Biology Group, International Centre for Genetic
Engineering and Biotechnology (ICGEB), Trieste, Italy.\\
$^{7}$ Bioinformatics Facility, Department of Cellular, Computational and
Integrative Biology (CIBIO), University of Trento, Trento, Italy.\\[0.5em]
\textsuperscript{$\dagger$}Moussa Kassim Mohamed and Martin Ocharo contributed equally.\\
\textsuperscript{*}Corresponding author: Stefano Cacciatore,
\href{mailto:stefano.cacciatore@icgeb.org}{stefano.cacciatore@icgeb.org}.

\section{Scope}

This supplement describes the datasets, software environment, calibration
grids, and validation procedures used in the article. It also provides
detailed comparison results and documents input handling, fitted-index
lifetime, device memory ownership, and experimental graph methods. The
accompanying replication files contain the parameter grids, recorded
outcomes, and checksums used to reconstruct the numerical summaries.

\section{Inputs, results, and index lifetime}

\subsection{Distances and numerical behavior} \label{sec:supp-numerics}

All returned Euclidean distances use the ordinary $L_2$ scale. For the
exact-reference checks, an observed distance $a$ and its reference value $b$
are treated as equal when
\[
|a-b|\leq 10^{-5}+10^{-4}|b|.
\]
The first term allows for float32 rounding near zero, and the second allows a
slightly larger difference for larger distances. These constants were fixed
before analysis and are used only to audit exact-method agreement; they do not
change search results or approximate-recall decisions. Native exact scorers
order ties by distance and then row identifier. Other libraries may return a
different row identifier when several observations are tied at the $k$th
distance. Such a result is accepted only when exact rescoring confirms that the
alternative neighbor has the same distance under the criterion above.

\subsection{Host-result and dispatch contract}

The function \code{nn()} returns a \class{faissR\_nn} list with
$m\times k$ matrices named \code{indices} and \code{distances}. Identifiers
are one-based R integers. Distances are doubles by default, or float32 with
\pkg{float} \citep{float}, ordered from smallest to largest. Their interpretation is
recorded in \code{distance\_is\_metric}, \code{distance\_semantics},
\code{distance\_comparable\_across\_queries}, and \code{distance\_order}.

The argument \code{backend} selects a device, while
\code{backend\_used} identifies an implementation, such as
\code{faiss\_hnsw}. The returned fields \code{requested\_method},
\code{requested\_backend}, and \code{resolved\_backend} distinguish the request
from its execution. Parameters and selection information are stored in nested
metadata. The application programming interface (API) retains these names
for compatibility.

On CPU, \code{exact}, \code{flat}, and \code{bruteforce} use FAISS Flat.
On CUDA, \code{exact} prefers FAISS GPU Flat and otherwise uses direct cuVS
brute force; \code{flat} requires FAISS GPU Flat; and \code{bruteforce}
prefers cuVS and otherwise uses FAISS GPU Flat. An unsupported explicit
method, metric, or device produces an error without substituting another.
Use \code{nn\_capabilities(runtime = TRUE)} to inspect supported combinations
and \code{backend\_info()} to inspect the linked libraries. This dispatch
uses capabilities provided by FAISS and cuVS
\citep{JohnsonDouzeJegou2021,DouzeEtAl2024,cuVS,FaissCuVS}.

\subsection{Fitted models}

Calling \code{knn()} without test data returns a
\class{faissR\_knn\_model}; with test data, it returns predictions directly.
For compatible CPU methods, \code{predict()} retains and reuses FAISS Flat,
hierarchical navigable small-world (HNSW), inverted file (IVF), or
product quantization (PQ) indexes in the IVF-PQ family
\citep{MalkovYashunin2020,JegouDouzeSchmid2011}. Reuse is reported as
\code{query\_source = "fitted\_index"}. Otherwise the same requested method
is rebuilt. Native allocations are released by finalizers when their
owning objects become unreachable; there is no manual public destroy function.

The secondary \code{fast\_kmeans()} interface returns a
\class{faissR\_kmeans} object with assignments, centers, within-cluster sums
of squares, sizes, convergence information, and execution metadata.
Its provider-backed implementation follows the usual k-means objective
\citep{MacQueen1967,Lloyd1982}. Predictive accuracy and clustering quality
were not compared in this study.

\subsection{GPU buffers and compiled consumers}

Calling \code{nn\_gpu()} returns a \class{faissR\_gpu\_knn} object with an
owning \code{handle} and non-owning \code{indices\_ptr} and
\code{distances\_ptr} views. The buffers are column-major $m\times k$ arrays
with one-based int32 identifiers and float32 distances. The object records
\code{result\_residency = "cuda"} and
\code{device\_to\_host\_result\_copies = 0}.
Both host and device identifiers are signed 32-bit integers, limiting the
interface to $2^{31}-1$ reference rows; practical limits are usually lower.

The producer synchronizes the recorded device before returning the object.
It exports neither a CUDA stream nor a FAISS/cuVS resource manager.
A consumer checks application binary interface (ABI) version 1, retains the
owning handle, selects the recorded device, and completes and synchronizes
its own work before releasing the owner. It must not free the child pointers.
A different resource manager can access these buffers after handoff, but
does not share the producer's stream, allocator, or ownership.

The handle's finalizer releases both CUDA allocations. Consumers must retain
it while asynchronous work is outstanding, rather than relying on garbage
collection for synchronization. The pointers are valid only in the same
process and on the same device. Serialization and CUDA interprocess
communication (IPC) are unsupported. Calling \code{gpu\_knn\_to\_host()}
copies both buffers to host matrices without invalidating the resident owner.

\section{Rationale and scope of the experimental graph refinements}
\label{sec:supp-derived}

These routes were introduced to cover gaps in the provider interfaces available
to the package. FAISS NSG supported only Euclidean search, and the linked FAISS
builds used during package development did not provide sufficiently reliable
NSG or NN-descent graph construction for the public interface. cuVS exposed
Vamana construction and serialization, but not a general Vamana query
interface. A package-owned self-search refinement path made it possible to use
the same bounded candidate representation on CPU and CUDA and to support
Euclidean, cosine, and correlation search after the metric transformations
described below \citep{DouzeEtAl2024,cuVS}.

The purpose was therefore to provide controlled fallbacks and to study selected
pruning or neighborhood-update operations, not to replace or improve the
published algorithms. The routes remain experimental, are excluded from the
principal performance comparisons, and are retained for reproducibility and
further method development. Their experimental status is reported in the
returned metadata, including when package-internal routing selects one of them.

The package-owned \code{nsg\_style}, \code{vamana\_style}, and CPU
\code{nndescent\_style} methods refine candidates for self
$k$-nearest-neighbor (kNN) search. They borrow selected ideas from navigating
spreading-out graph (NSG), Vamana, and nearest-neighbor descent (NN-descent),
but they do not reproduce the published algorithms
\citep{FuEtAl2019,SubramanyaEtAl2019,DongMosesLi2011}. Their results set
\code{implementation\_status = "experimental"}, \code{experimental = TRUE},
and \code{canonical\_reimplementation = FALSE}, with scope
\code{package\_owned\_style\_implementation}. CUDA \code{nndescent\_style}
may instead call cuVS NN-descent directly; metadata then identify an
external-provider route.

These methods differ from the published algorithms in their algorithmic
structure, not only in their tuning parameters. Published NSG and Vamana
construct navigable graph indexes and specify graph traversal for later
queries. Published NN-descent constructs an approximate kNN graph through a
particular iterative update procedure. The package-owned methods instead
operate on a bounded candidate set for each row of a self-search. They retain
one pruning or neighborhood-update operation and compute exact distances
within the resulting candidate set. They do not produce the canonical graph
index or implement its complete construction and query procedure. The
\code{\_style} suffix denotes this restricted relationship.

Cosine input is row-normalized; correlation input is row-centered and
row-normalized. Both are refined as Euclidean candidates before the
documented distances are restored. The following descriptions refer to
package-owned stages, not to the full canonical algorithms.

\paragraph{Published NSG and \code{nsg\_style}.}
Published NSG starts from an approximate kNN graph, chooses a navigating node
near the data centroid, gathers candidates by graph search, applies a
monotonic-relative-neighborhood-graph pruning rule, and repairs connectivity.
Queries subsequently traverse the resulting graph from the navigating node.
The package-owned route does not perform these index-construction or query
stages. Instead, given an ordered seed-neighbor matrix $S$ with width $s$,
output degree $r$, and protected prefix $k$, it proceeds as follows:
\begin{enumerate}
\item For each query row $q$, remove duplicates and the self identifier from
      $S_q$, preserving order, and retain the first $k$ candidates.
\item For every later candidate $v$, reject $v$ when a retained $u$ satisfies
      $d(u,v)<d(q,v)$; otherwise retain $v$, up to degree $r$.
\item Fill unoccupied degree slots from the original ordered candidates and
      exactly rescore the retained row candidates; return the closest $k$.
\end{enumerate}
Thus, \code{nsg\_style} retains one NSG-like relative-neighborhood test but
does not construct an NSG index, select a navigating node, repair graph
connectivity, or implement NSG query traversal.

\paragraph{Published Vamana and \code{vamana\_style}.}
Published Vamana initializes a random bounded-degree directed graph, selects a
medoid, processes vertices in random order, and alternates greedy graph search
with robust pruning. Retained outgoing edges are inserted reciprocally, and
overfull neighbor lists are pruned again. Construction uses two passes, first
with $\alpha=1$ and then with the requested $\alpha$. The package-owned route
does not implement that incremental construction. It uses the same ordered
seed and exact row-rescoring framework as \code{nsg\_style}, but rejects a
candidate $v$ when a retained $u$ satisfies
$\alpha d(u,v)\leq d(q,v)$. It therefore retains the robust-pruning inequality
but omits random graph initialization, medoid-based greedy search, random-order
insertion, reciprocal edge updates and repruning, the two construction passes,
and traversal of a reusable Vamana index. It also does not implement the
solid-state-drive layout and compression system in DiskANN, within which
Vamana was introduced.

\paragraph{Published NN-descent and \code{nndescent\_style}.}
Published NN-descent initializes an approximate kNN graph randomly and improves
it by repeatedly joining new and old neighbors, including sampled reverse
neighbors. Its update count supplies a convergence-based stopping rule. The
package-owned CPU route changes both initialization and iteration. For pool
size $L$, candidate cap $C$, projection count $P$, fixed iteration count $t$,
and integer seed, it proceeds as follows:
\begin{enumerate}
\item Generate $P$ seeded Gaussian directions, sort projected rows with an
      identifier tie break, and collect neighbors from bounded projection
      windows. Seeded random and then sequential fill complete short pools.
\item At every fixed iteration, freeze the previous graph, construct bounded
      reverse-neighbor lists, and evaluate capped neighbor-of-neighbor and
      reverse-neighbor joins for each row.
\item Retain the best $L$ candidates after each iteration and return the first
      $k$ after the final exact candidate scoring.
\end{enumerate}
The implementation therefore uses projection-based rather than purely random
initialization, bounded reverse lists and candidate joins, and a fixed number
of synchronous iterations. It omits the published new/old flag schedule,
sampled local joins, and update-based convergence rule. It returns the result
of a self-search rather than a general NN-descent query index. For fixed input,
seed, parameters, and seed provider, package-owned stages use deterministic
ordering. An external seed provider can retain its documented floating-point
or threading variation.
Table~\ref{tab:supp-derived-methods} summarizes computational bounds and the
scope of the evidence for these refinements and other secondary routes.

\begin{table}[H]
\centering

\small
\begin{tabularx}{\textwidth}{@{}>{\raggedright\arraybackslash}p{0.19\textwidth}>{\raggedright\arraybackslash}p{0.43\textwidth}X@{}}
\toprule
Route & Relationship to the published algorithm & Bounds and evidence in this article \\
\midrule
\code{nsg\_style} & Retains an NSG-like relative-neighborhood pruning test on
an existing seed list. It omits navigating-node selection, search-based
candidate collection, connectivity repair, and canonical graph traversal. &
Exact seeding can dominate at $O(n^2p)$; pruning is at most $O(nr^2p)$ and
rescoring $O(nrp)$. Experimental; excluded from principal comparisons. \\
\code{vamana\_style} & Retains Vamana's $\alpha$-scaled robust-pruning test on
an existing seed list. It omits random initialization, incremental two-pass
construction, medoid search, reciprocal edge insertion and repruning, and
canonical graph traversal. & Same package-route bounds as
\code{nsg\_style}. Experimental; excluded from principal comparisons. \\
CPU \code{nndescent\_style} & Retains bounded neighbor-of-neighbor and
reverse-neighbor updates. Projection seeding and fixed synchronous iterations
replace random initialization, new/old sampled joins, and convergence-based
stopping. & Initialization is approximately $O(Pnp+Pn\log n+nLp)$; refinement
is bounded by $O(tnCp)$. Experimental; excluded from principal comparisons. \\
Grid, $d\in\{2,3\}$ & Package-owned exact spatial-cell search, rather than a
variant of the three graph algorithms above. & Build and memory $O(n+B^d)$;
occupancy-dependent query, worst-case $O(n)$ per query. Reference-audited. \\
IVF-PQ/FastScan \citep{FaissFastScan} & Provider-defined FAISS/cuVS
quantization and search, rather than a package-owned graph refinement. &
Supported provider interfaces; no route-specific principal performance claim. \\
\bottomrule
\end{tabularx}
\caption{Relationship between package-owned secondary routes and the
published algorithms, with implementation-level bounds and evidence scope.
$n$ is row count, $p$ dimension, $s$ seed width, $r$ retained degree, $L$ pool
size, $C$ candidate cap, $P$ projection count, $t$ iterations, and $B$ cells
per grid axis. The bounds apply to the package-owned implementations.}
\label{tab:supp-derived-methods}
\end{table}

Automatic method and parameter lookups are deterministic. Package-owned graph stages obtain their
seed from the R option \path{faissR.approx_knn_seed}. Results record the seed,
seed provider, and parameters. FAISS and cuVS construction can include random
or parallel stages; GPU execution is not promised to be bitwise reproducible
across drivers, devices, or provider versions. Reproduction therefore
requires the input fingerprint, source and experiment commits, requested
method and backend, resolved provider, metric, $k$, recall tier, parameters,
seed, thread controls, software versions, and hardware identity.

\section{Installation and portability}
\label{sec:supp-installation}

The installation requirements and discovery rules are given in the main
article. Table~\ref{tab:install-matrix} records the environments in which those
paths were tested and distinguishes functional from diagnostic-only builds.

\begin{table}[H]
\centering

\small
\begin{tabularx}{\textwidth}{@{}p{0.19\textwidth}p{0.25\textwidth}X@{}}
\toprule
Status & Platform and providers & Evidence or boundary \\
\midrule
Functional, completed & macOS arm64; R 4.6.0; FAISS 1.14.3; libomp 22.1.x; Homebrew clang 22.1.1; gfortran 12.2.0 & Strict source install and CPU search smoke passed. The retained package-check log reports no errors or warnings; manual rendering is checked separately. \\
Functional, completed & Debian 13; R 4.5.3; FAISS 1.14.3; cuVS 26.06; CUDA 13.2; NVIDIA L40S, driver 595.58.03 & Publication route-contract and benchmark environment. The compiler executable/version used to construct the image was not retained and is therefore not inferred. \\
Diagnostic only & Recognized Windows, macOS staging, or WebAssembly builders without compatible FAISS & Package can load and report the missing provider; nearest-neighbor calls are unavailable. This is portability diagnostics, not functional support. \\
\bottomrule
\end{tabularx}
\caption{Recorded installation environments and diagnostic-build boundaries.}
\label{tab:install-matrix}
\end{table}

The installation vignette gives command sequences for a Linux CPU-only
FAISS 1.14.3 source installation, a Homebrew macOS CPU build, and a strict
Linux CUDA/cuVS build. It includes capability checks and troubleshooting for
missing FAISS, \code{omp.h}, runtime \code{libfaiss},
\code{GLIBCXX} symbol-version mismatches, \code{nvcc}, and cuVS.
External-library upgrades require a fresh compatibility check. The repository
also provides an isolated Valgrind test of exact and HNSW calls to detect
common native-memory errors.

\pkg{faissR} and FAISS use the Expat License; R package metadata records this
template with the identifier \code{MIT + file LICENSE}. cuVS uses Apache-2.0.
The optional CUDA toolkit and driver are governed by NVIDIA's software development kit
(SDK) terms. These native dependencies are supplied by the user or system,
not copied into the package source. The installed
\path{THIRD_PARTY_LICENSES.md} records the relevant links and distribution
boundaries.

\section{Calibration}

\subsection{Parameter grids and selection}

The main article defines target recall and the selection rule. Here,
Table~\ref{tab:grid} documents the parameter grids crossed with device,
distance, $k\in\{15,30,50,100\}$, and the three requested recall tiers.

\begin{table}[h]
\centering

\begin{tabularx}{\textwidth}{@{}lY@{}}
\toprule
Method & Candidate values or construction rule \\
\midrule
Exhaustive Flat/brute force & CPU batches $\{2048,8192,32768,65536,131072\}$ and cache off/on; CUDA batches $\{1024,4096,8192,16384,32768,65536\}$ and resource reuse on/off. \\
Grid & Shape-dependent cell width and expansion, restricted to supported low-dimensional Euclidean cells. \\
HNSW & CPU $M\in\{6,8,10,12,16,24,32,48,64\}$ with paired construction/search schedules; CUDA uses nine paired graph/intermediate/search-width schedules. \\
IVF & Shape- and $k$-dependent base \code{nlist}/\code{nprobe}, crossed with nine prespecified multiplier pairs from coarse probing to full probing. \\
IVF-PQ & IVF controls, valid divisors near $\{16,32,48,64,96\}$ on CPU or $\{16,32,48,64\}$ on CUDA, 8-bit codes, and query batch. \\
IVF-PQ FastScan & Valid product-quantizer divisors, 4-bit codes, block size $\{32,64\}$ on CPU, and a nine-level refinement schedule. \\
CAGRA & Nine graph/intermediate/search schedules crossed with \code{auto}, IVF-PQ, NN-descent, and iterative-CAGRA build strategies where supported. \\
faissR NN-descent-derived refinement & Prespecified candidate-pool, graph-degree, and iteration schedules; CUDA schedules are shape- and $k$-dependent. \\
faissR NSG-/Vamana-derived refinements & Output degree $\{4,6,8,12,16,24,32,48,64,96\}$ with paired seed/search breadth; Vamana also uses $\alpha\in[1,1.5]$. \\
\bottomrule
\end{tabularx}
\caption{Method-specific calibration grid summary. The complete
machine-readable grid and its source-grid checksums are available in the
\href{https://github.com/tkcaccia/faissR/tree/main/manuscript/jss}{public
GitHub replication materials}.}
\label{tab:grid}
\end{table}

The complete machine-readable grid contains 41,229 configurations for the
three metrics, generated from 63 archived source grids. The accompanying
manifest records their checksums. These materials specify the actual values
for every dataset, device, method, and $k$, including adjustments needed to
satisfy library constraints such as a valid product-quantizer dimension.

\subsection{Local HNSW pre-tuning protocol} \label{sec:local-hnsw}

The optional \code{hnsw\_tune()} workflow is separate from the bundled
shape-based policy and from the publication calibration. Its initial scope is
CPU FAISS HNSW with Euclidean distance. Sparse or delayed inputs, unsupported
metrics, and unavailable FAISS support produce errors; the procedure does not
substitute another method or device.

The procedure samples at most 20,000 reference rows without using a prefix of
the input. It then draws disjoint tuning and confirmation query sets from that
reference. Exact FAISS Flat search supplies the neighbor reference. Recall is
tie-aware: an alternative identifier at the $k$th boundary is accepted when
its distance is within $10^{-5}+10^{-4}|d_k|$ of the exact boundary distance.
The caller prespecifies query rows per batch and the expected number of batches
served by one index. For each candidate pair of $M$ and construction effort,
graph construction is repeated and the final graph is queried repeatedly at
each candidate search effort on the same host. The output reports mean,
fifth-percentile, tenth-percentile, and minimum query recall, together with an
empirical 95\% query-bootstrap lower bound.

Eligibility on the tuning queries requires the lower bound to reach the
requested tier plus a prespecified safety margin, capped at one. Among eligible
candidates, selection minimizes the projected workload time. This objective is
the median pilot construction time, scaled by the full-to-pilot reference-row
ratio, plus median query time scaled by the requested-to-timed query-row ratio
and multiplied by the expected number of batches. Median query time,
construction time, and then $(M,$ construction effort, search effort$)$ provide
deterministic tie breaking. This is an explicit linear projection from a
bounded pilot, not a claim that HNSW construction scales exactly linearly. The
selected setting must independently reach the requested
lower bound on the untouched confirmation queries before a full-data index is
built. The returned object records sampled row identifiers, selected settings,
the workload objective, timing dispersion, lower-tail recall, exact-reference
and pilot timing, provider information, and whether confirmation succeeded.
Its owned index pointer can be queried
repeatedly with \code{predict()} and is valid only within the creating R
session.

The disjoint query split evaluates within-dataset confirmation. It does not
test whether a rule developed on one dataset transfers to another dataset with
similar dimensions but different neighborhood geometry. Unless the complete
reference is used, it also estimates recall on the sampled reference rather
than proving recall against the full collection.

The default sweep is bounded by the number of graph builds, elapsed pilot
time, and an optional estimate-based memory limit. The time limit is a soft
cutoff checked between graph candidates; an active build or search is not
interrupted. If the sweep stops early, no deployable setting is returned. For
comparative performance, adaptive end-to-end time comprises
exact-reference generation, all pilot builds and searches, final index
construction, and the application query. Fitted-index query time is a separate
quantity, and no speed ratio is defined when the requested recall criterion is
not met.

\subsection{Completion and timing sensitivity}

Table~\ref{tab:calibration-missing} disaggregates the unavailable operating
points reported in the main article by device, method, and failure reason.

\begin{table}[h]
\centering

\begin{tabular}{@{}llrr@{}}
\toprule
Backend & Method & Time limit & Memory kill \\
\midrule
CPU & Brute force & 72 & 0 \\
CPU & Exact & 72 & 0 \\
CPU & Flat & 72 & 0 \\
CPU & IVF & 12 & 12 \\
CPU & NSG-derived refinement & 42 & 0 \\
CPU & Vamana-derived refinement & 36 & 0 \\
CUDA & NN-descent-derived refinement & 0 & 9 \\
CUDA & NSG-derived refinement & 0 & 12 \\
CUDA & Vamana-derived refinement & 0 & 12 \\
\midrule
Total & & 306 & 45 \\
\bottomrule
\end{tabular}
\caption{Unavailable calibration operating points by backend, method, and
recorded failure class.}
\label{tab:calibration-missing}
\end{table}

The unavailable totals were 117 for correlation, 123 for cosine, and 111 for
Euclidean distance. By dataset, they were 201 for FR-FCM-ZYRM, 144 for ImageNet
features, and 6 for flow18. The machine-readable results cross device,
method, metric, dataset, and failure reason.
Table~\ref{tab:calibration-stability} summarizes the timing gaps between
the fastest and second-fastest eligible candidates.

\begin{table}[h]
\centering

\begin{tabular}{@{}lrrr@{}}
\toprule
Comparison & Points & Median margin & Within 5\% \\
\midrule
Configuration, all & 5,285 & 2.67\% & 60.8\% \\
Configuration, CPU & 2,462 & 6.71\% & 43.6\% \\
Configuration, CUDA & 2,823 & 0.73\% & 75.8\% \\
Method family, all & 641 & 50.24\% & 10.1\% \\
\bottomrule
\end{tabular}
\caption{Selection sensitivity from eligible-candidate timing margins. The
margin is $T_{(2)}/T_{(1)}-1$, where $T_{(1)}$ and $T_{(2)}$ are the fastest and
second-fastest eligible times. These are near-tie diagnostics, not repeated
timing estimates.}
\label{tab:calibration-stability}
\end{table}

With only one successful timing per configuration, the data cannot establish
how often the selected configuration would change across timing replicates.
The gap between the fastest and second-fastest eligible times instead shows
how sensitive the choice could be to modest timing variation.
Configuration choices had the smallest gaps on CUDA; different method
families were generally better separated.
Table~\ref{tab:calibration-validation} compares calibration with validation.

\begin{table}[h]
\centering

\begin{tabular}{@{}llrrr@{}}
\toprule
Backend & Route & Points & Recall change & Time ratio \\
\midrule
CPU & HNSW & 324 & $+0.0053$ & 1.05 \\
CUDA & Automatic method, all & 324 & $0.0000$ & 0.96 \\
CUDA & Automatic method, Flat & 280 & $0.0000$ & 0.95 \\
CUDA & Automatic method, IVF & 44 & $+0.0096$ & 2.12 \\
\bottomrule
\end{tabular}
\caption{Calibration-to-validation change. Recall change is validation mean
recall minus calibration recall; time ratio is median validation time divided by
calibration time.}
\label{tab:calibration-validation}
\end{table}

\section{Datasets and exact references}

Table~\ref{tab:dataset-provenance} records the source, preprocessing,
redistribution status, and final representation for every benchmark matrix.
The accompanying machine-readable provenance file additionally supplies
release identifiers, acquisition instructions, terms, and fingerprints.

\begin{table}[H]
\centering
\scriptsize

\setlength{\tabcolsep}{3pt}
\begin{tabularx}{\textwidth}{@{}p{0.13\textwidth}p{0.23\textwidth}X p{0.09\textwidth}@{}}
\toprule
Dataset & Source or accession & Final input and preprocessing & Converted matrix redistributed \\
\midrule
COIL20 & Columbia processed COIL-20 archive, \code{coil-20-proc.zip} \citep{Nene1996} &
1,440 $\times$ 16,384; processed PNG images converted to grayscale when needed and flattened; no scaling, PCA, or other reduction. & No \\
FashionMNIST & Official Fashion-MNIST training and test IDX release \citep{Xiao2017} &
70,000 $\times$ 784; official splits concatenated and $28\times28$ grayscale pixels flattened; no additional scaling or reduction. & Yes, MIT \\
flow18 & opt-SNE Figshare item 9927986, version 1 \citep{Belkina2019} &
1,000,021 $\times$ 11; published matrix values used without further normalization, filtering, or transformation. & Yes, CC BY 4.0 \\
FR-FCM-ZYRM & FlowRepository accession \code{FR-FCM-ZYRM} \citep{VanUnen2017} &
5,220,347 $\times$ 32; 102 FCS files, 32 biological-marker channels, $\operatorname{asinh}(x/5)$ transformation. & No \\
ImageNet features & Authorized ILSVRC 2012 training images represented with DINOv2 \citep{Russakovsky2015,Oquab2024} &
1,281,167 $\times$ 1,024; precomputed DINOv2 embeddings and class indices, with no benchmark-time scaling or reduction. & No \\
mass41 & opt-SNE Figshare item 9927986, version 1 \citep{Belkina2019} &
965,282 $\times$ 14; published matrix values, already $\operatorname{asinh}(x/5)$ transformed, used without further processing. & Yes, CC BY 4.0 \\
MetRef & \pkg{KODAMA} source release, dataset \code{MetRef} \citep{Cacciatore2014} &
873 $\times$ 375; zero-column-sum variables removed, followed by KODAMA normalization and scaling; no PCA or embedding. & No \\
MNIST & Training and test IDX files from the recorded public mirror \citep{LeCun1998} &
70,000 $\times$ 784; official splits concatenated and $28\times28$ pixels flattened; no additional scaling or reduction. & No \\
USPS & \pkg{Rdimtools} source release, dataset \code{usps}
\citep{YouShung2022,Hull1994} &
11,000 $\times$ 256; package feature matrix retained without added scaling, PCA, or dimensionality reduction. & No \\
\bottomrule
\end{tabularx}
\caption{Dataset provenance and benchmark representation. ``No'' means that
the converted matrix is not included in the replication archive; users obtain
the named source and run the tracked conversion and fingerprint checks.}
\label{tab:dataset-provenance}
\end{table}

Converted inputs are distributed only when the source terms clearly permit
redistribution. In particular, reproducing ImageNet features requires
authorized access to the original images and regeneration of the recorded
DINOv2 representation. Dataset terms are separate from the package license.
The acquisition code in \path{benchmark_scripts/jss_reproduction/} checks
fingerprints and records missing or nonmatching inputs as unavailable.

Table~\ref{tab:reference-key} decomposes the 87-record reference archive
reported in the main article. The row-level key stores dimensions, seeds,
fingerprints, providers, and independent CPU cross-audit statistics.

\begin{table}[h]
\centering

\begin{tabular}{@{}lrrrrr@{}}
\toprule
Group & Objects & Metrics & Seeds & Queries per record & Records \\
\midrule
Real datasets & 9 & 3 & 3 & $\min(n,1024)$ & 81 \\
Synthetic spatial shapes & 2 & 1 (Euclidean) & 3 & 1024 & 6 \\
\midrule
Total & 11 & -- & -- & -- & 87 \\
\bottomrule
\end{tabular}
\caption{Exact-reference archive key. The complete machine-readable key,
including dimensions, seeds, fingerprints, providers, and CPU cross-audit
statistics, is available in the
\href{https://github.com/tkcaccia/faissR/tree/main/manuscript/jss}{public
GitHub replication materials}.}
\label{tab:reference-key}
\end{table}

\section{Computing environment}

Table~\ref{tab:supp-environment} reports the hardware used for the experiments.
The CPU jobs ran on several generations of x86-64 processor, so a single model
cannot represent all runs. Paired CPU comparisons used the same node and
allocation, and CUDA timers ended after device synchronization.

\begin{table}[h]
\centering

\small
\begin{tabularx}{\textwidth}{@{}>{\raggedright\arraybackslash}p{0.30\textwidth}Y@{}}
\toprule
Component & Hardware configuration \\
\midrule
CPU nodes & Heterogeneous x86-64 compute nodes spanning several processor
generations \\
CPU parallelism & 12 requested threads per job \\
GPU & NVIDIA L40S with 46,068 mebibytes (MiB) of device memory and compute
capability 8.9 \\
\bottomrule
\end{tabularx}
\caption{Benchmark hardware.}
\label{tab:supp-environment}
\end{table}

The jobs ran under Debian 13 (trixie) Linux in a Singularity container
\citep{Singularity}, with \proglang{R} 4.5.3
\citep{R}, OpenBLAS 0.3.33, LAPACK 3.12.0, FAISS 1.14.3,
\code{libcuvs} 26.06, CUDA 13.2, and NVIDIA driver 595.58.03. The external CPU
comparators were \pkg{BiocNeighbors} 2.4.0, \pkg{FNN} 1.1.4.1,
\pkg{RANN} 2.6.2, \pkg{rnndescent} 0.2.0, \pkg{Rnanoflann} 0.0.3,
\pkg{RcppAnnoy} 0.0.23, and \pkg{RcppHNSW} 0.7.0. The frozen container had
SHA-256 digest
\texttt{0cd4d0df406bd007}\allowbreak\texttt{5046b16d2e8a4d3a}\allowbreak
\texttt{e78ee61d98e6d47c}\allowbreak\texttt{639986e28ea6f203}.

Calibration, reference, and independent-query timing used direct float32 input.
The controlled CPU comparison instead supplied the same \proglang{R} double
matrix to both routes, so the \pkg{faissR} conversion to float32 remained inside
its timer. Paired ratios were summarized within datasets before datasets were
combined.

\section{CPU comparisons}
\label{sec:supp-cpu}

This section supplies the provider settings and dataset-level results behind
the CPU comparisons in the main article. The external interfaces are
documented by \citet{BiocNeighbors}, \citet{FNN}, \citet{RANN},
\citet{rnndescent}, \citet{Rnanoflann}, \citet{RcppAnnoy}, and
\citet{RcppHNSW}.

\subsection{HNSW at fixed settings}

Each Slurm task fixed the dataset, $k$, and comparator. \pkg{faissR} and either
\pkg{BiocNeighbors} or \pkg{RcppHNSW} ran on the same node and allocation.
The first route was randomized reproducibly within each validation seed, then
the order alternated over five repetitions. Each repetition used an isolated
worker and an untimed 128-row warm-up. Both interfaces received the same R
double matrix.

Comparator settings were chosen before execution. \pkg{BiocNeighbors} used
\code{nlinks = 16}, \code{ef.construction = 200}, and
\code{ef.search = max(50, 3*k)}. The corresponding \pkg{RcppHNSW} settings
were \code{M = 16}, \code{ef\_construction = 200}, and
\code{ef = max(50, 3*k)}. These settings were not optimized separately for
each provider. A pair entered the timing summary only if both methods
achieved mean recall of at least 0.99 in every included repeat and the audit
confirmed a common hostname, allocation, and opposing execution order.

Table~\ref{tab:supp-cpu-paired} and Figure~\ref{fig:supp-paired-cpu} report
the fixed-configuration cold-call results by dataset. Build and fitted-query
ratios are omitted because this experiment did not establish equivalent
fitted-query recall for both methods.

\begin{table}[ht]
\centering
\small

\begin{tabular}{@{}lrrr|rrr@{}}
\toprule
& \multicolumn{3}{c}{\pkg{BiocNeighbors}} & \multicolumn{3}{c}{\pkg{RcppHNSW}} \\
Dataset & Complete & Target & Ratio & Complete & Target & Ratio \\
\midrule
COIL20 & 40 & 40 & 15.27 & 40 & 40 & 2.70 \\
FashionMNIST & 40 & 40 & 11.40 & 40 & 40 & 2.22 \\
FR-FCM-ZYRM & 21 & 0 & -- & 40 & 10 & 0.81 \\
flow18 & 40 & 40 & 2.83 & 40 & 40 & 0.72 \\
ImageNet features & 30 & 0 & -- & 40 & 0 & -- \\
mass41 & 40 & 40 & 2.72 & 40 & 40 & 0.70 \\
MetRef & 40 & 40 & 4.79 & 40 & 40 & 1.82 \\
MNIST & 40 & 40 & 11.57 & 40 & 40 & 1.87 \\
USPS & 40 & 13 & 23.66 & 40 & 18 & 4.90 \\
\bottomrule
\end{tabular}
\caption{Controlled CPU HNSW cold-call results by dataset. Entries give
completed pairs, mean-recall-matched pairs, and the median
$T_{\mathrm{comparator}}/T_{\mathrm{faissR}}$ among mean-recall-matched pairs.
A dash means that no completed pair attained 0.99 for both routes.}
\label{tab:supp-cpu-paired}
\end{table}

\begin{figure}[ht]
\centering
\includegraphics[width=\textwidth]{fig_paired_cpu_log_ratio.pdf}
\caption{Controlled CPU HNSW cold-call results for the prespecified public
configurations. Points are dataset medians and horizontal lines are
interquartile intervals among pairs for which both routes had mean recall of
at least 0.99 in every included replicate. Ratios are
$T_{\mathrm{comparator}}/T_{\mathrm{faissR}}$ on a logarithmic axis. Black
circles and grey triangles distinguish the comparators in grayscale; crosses
indicate that no pair met the recall criterion. The right-hand columns report
eligible pairs (\code{n}) and timeouts (\code{t}) among 40 planned pairs per
dataset and comparator.}
\label{fig:supp-paired-cpu}
\end{figure}
\clearpage

\subsection{Seven-package comparison}

Table~\ref{tab:supp-comprehensive-datasets} and
Figure~\ref{fig:supp-comprehensive-r} expand the comparison classes in the
main article into dataset-level results. Panel D additionally reports the
experimental \pkg{rnndescent} comparison, which is excluded from the principal
performance table.
\begin{table}[H]
\centering
\scriptsize
\begin{tabularx}{\textwidth}{@{}l*{4}{>{\centering\arraybackslash}X}@{}}
\toprule
\multicolumn{5}{@{}l}{\textit{Panel A: same-family exhaustive comparisons}} \\
\addlinespace[0.25em]
Dataset & \pkg{FNN} exact & \pkg{RANN} exact & \pkg{Rnanoflann} exact & \pkg{BiocNeighbors} exact \\
\midrule
COIL-20 & 32.00 (72) & 37.69 (48) & 8.60 (24) & 4.73 (48) \\
Fashion-MNIST & -- & -- & 22.04 (24) & 15.33 (48) \\
FR-FCM-ZYRM & -- & -- & -- & -- \\
MNIST & -- & -- & 21.35 (24) & 15.40 (48) \\
MetRef & 4.24 (72) & 6.63 (48) & 2.31 (24) & 2.21 (48) \\
USPS & 37.80 (45) & 40.73 (18) & 9.00 (18) & 3.16 (36) \\
flow18 & -- & -- & -- & 2.06 (6) \\
ImageNet features & -- & -- & -- & -- \\
mass41 & -- & -- & -- & 1.97 (4) \\
\midrule
\multicolumn{5}{@{}l}{\textit{Panel B: same-family HNSW comparisons}} \\
\addlinespace[0.25em]
Dataset & \multicolumn{2}{c}{\pkg{BiocNeighbors} HNSW} & \multicolumn{2}{c}{\pkg{RcppHNSW} HNSW} \\
\midrule
COIL-20 & \multicolumn{2}{c}{11.58 (48)} & \multicolumn{2}{c}{2.61 (48)} \\
Fashion-MNIST & \multicolumn{2}{c}{9.84 (42)} & \multicolumn{2}{c}{1.67 (48)} \\
FR-FCM-ZYRM & \multicolumn{2}{c}{--} & \multicolumn{2}{c}{0.75 (12)} \\
MNIST & \multicolumn{2}{c}{8.88 (48)} & \multicolumn{2}{c}{1.68 (48)} \\
MetRef & \multicolumn{2}{c}{5.19 (48)} & \multicolumn{2}{c}{6.60 (48)} \\
USPS & \multicolumn{2}{c}{13.36 (17)} & \multicolumn{2}{c}{3.95 (17)} \\
flow18 & \multicolumn{2}{c}{2.98 (48)} & \multicolumn{2}{c}{0.78 (48)} \\
ImageNet features & \multicolumn{2}{c}{--} & \multicolumn{2}{c}{--} \\
mass41 & \multicolumn{2}{c}{2.78 (48)} & \multicolumn{2}{c}{0.73 (48)} \\
\midrule
\multicolumn{5}{@{}l}{\textit{Panel C: task-level Annoy alternatives}} \\
\addlinespace[0.25em]
Dataset & \multicolumn{2}{c}{\pkg{BiocNeighbors} Annoy / \pkg{faissR} auto} & \multicolumn{2}{c}{\pkg{RcppAnnoy} / \pkg{faissR} auto} \\
\midrule
COIL-20 & \multicolumn{2}{c}{1.57 (48)} & \multicolumn{2}{c}{1.68 (48)} \\
flow18 & \multicolumn{2}{c}{1.89 (45)} & \multicolumn{2}{c}{2.62 (41)} \\
mass41 & \multicolumn{2}{c}{1.36 (3)} & \multicolumn{2}{c}{--} \\
\midrule
\multicolumn{5}{@{}l}{\textit{Panel D: experimental derived-route comparison}} \\
\addlinespace[0.25em]
Dataset & \multicolumn{4}{c}{\pkg{rnndescent} NN-descent / \pkg{faissR} \code{nndescent\_style}} \\
\midrule
COIL-20 & \multicolumn{4}{c}{0.22 (72)} \\
MetRef & \multicolumn{4}{c}{1.30 (72)} \\
USPS & \multicolumn{4}{c}{0.73 (36)} \\
\bottomrule
\end{tabularx}
\caption{Dataset medians for the seven-package comparison. Entries are median
$T_{\mathrm{comparator}}/T_{\mathrm{faissR}}$ followed by the number of
eligible paired repetitions in parentheses. A dash indicates that the dataset
had no pair in which both methods completed and met the requested mean recall.
Panels A and B are same-family comparisons. Panel C compares task-level
alternatives, not Annoy implementations. Panel D is an experimental
derived-route comparison and is excluded from the principal performance
table.}
\label{tab:supp-comprehensive-datasets}
\end{table}

\begin{figure}[h]
\centering
\includegraphics[width=0.94\textwidth]{fig_comprehensive_r_log_ratio.pdf}
\caption{Matched-node public-interface CPU comparisons. Open circles are
medians of eligible repetition-level ratios within each dataset. Solid points
and bars show the median and interquartile interval of those dataset medians.
Right-hand labels give matched-pair counts (\code{n}), datasets with at least
one matched pair (\code{d}), and timeouts (\code{t}) in each comparison class.
Each vertical-axis label names the external package and method followed by the
matched \pkg{faissR} method. Task-level and experimental comparisons are
identified in the labels.}
\label{fig:supp-comprehensive-r}
\end{figure}
\FloatBarrier

\subsection{Independently tuned HNSW}

The main article gives the tuning design and phase-specific results. The
following details clarify representation, self-neighbor handling, and thread
instrumentation for that experiment.

Both interfaces received the same R double matrix. Conversion to the float32
representation required by FAISS was included in the timed \pkg{faissR}
call; the comparators' internal precision was not instrumented. Calls
requested $k+1$ neighbors and removed the known source row within the timer:
\pkg{faissR} did so in compiled code, while comparator wrappers did so in R.
Provider-returned ordering was retained. The timings cover each complete
public interface, including its result handling.

Twelve threads were requested through the public interfaces. Thread controls
included Open Multi-Processing (OpenMP), OpenBLAS, and the Intel Math Kernel
Library (MKL), through \code{OMP\_NUM\_THREADS},
\code{OPENBLAS\_NUM\_THREADS}, and \code{MKL\_NUM\_THREADS}. The run recorded
the Slurm allocation, Linux affinity mask, these variables, and process-thread
counts before and after each call. They document requested limits and observed
process counts; active internal workers may still differ among providers.

\section{Query batches and index reuse}
\label{sec:supp-workload}

The query-workload experiment used five datasets and external batches of 1,
32, and up to 1,024 queries. It recorded cold and warm calls after clearing
the package cache at the start of each sequence. The difference
$\max(T_{\mathrm{cold}}-T_{\mathrm{warm}},0)$ estimates setup cost; it is not
a separately instrumented index-build time. For $A$ batches, the estimated
total was $T_{\mathrm{cold}}+(A-1)T_{\mathrm{warm}}$, evaluated at
$A\in\{1,10,100\}$.

The break-even point reported in the main article is the smallest $A$
satisfying
\[
T_{\mathrm{cold},M}+(A-1)T_{\mathrm{warm},M}
\leq A\,T_{\mathrm{cold},\mathrm{Flat}}.
\]
The Flat baseline rebuilt its index for every call; the comparison therefore
does not represent a reusable exact index.

GPU-resident output was assessed through interface, ownership, and explicit
transfer tests. Transfer speed, downstream consumer performance, and
quantitative host- or device-memory comparisons remain unmeasured.
Process-lifetime memory high-water marks cannot establish per-cell peak memory
when several cells share a worker.

\section{CUDA selection and sensitivity analyses}
\label{sec:supp-selector}

\subsection{Validation and route choices}

By metric, Flat/IVF counts were 64/44 for Euclidean, 108/0 for cosine, and
108/0 for correlation. By target they were 92/16 at 0.90, 92/16 at 0.95,
and 96/12 at 0.99. Each value of $k$ contributed 81 cells, with 70 Flat and
11 IVF selections. Table~\ref{tab:supp-auto-routes} gives the dataset counts.

\begin{table}[h]
\centering

\begin{tabularx}{\textwidth}{@{}Xrr@{}}
\toprule
Dataset & Flat & IVF \\
\midrule
COIL20 & 36 & 0 \\
FashionMNIST & 36 & 0 \\
FR-FCM-ZYRM & 24 & 12 \\
flow18 & 24 & 12 \\
ImageNet features & 28 & 8 \\
mass41 & 24 & 12 \\
MetRef & 36 & 0 \\
MNIST & 36 & 0 \\
USPS & 36 & 0 \\
\bottomrule
\end{tabularx}
\caption{CUDA automatic method selections by dataset. Each row contains 36
metric--$k$--target cells.}
\label{tab:supp-auto-routes}
\end{table}

The selected IVF parameters are listed in
Table~\ref{tab:auto-ivf-parameters}. Across 108 dataset--metric--$k$
combinations, changing the requested tier changed the method in four,
changed parameters without changing method in six, and changed neither in
98. The effect was concentrated in large Euclidean problems. For the
low-dimensional datasets, numeric parameters changed at $k=15$ and $k=30$,
but not at $k=50$ or $k=100$. ImageNet switched to Flat at target 0.99.

Probe counts obey $1\leq\code{nprobe}\leq\code{nlist}$. For example,
308 of 5,898 lists is about 5.2\%, and 160 of 1,590 is about 10.1\%.
These settings were feasible on the 48-gigabyte L40S. Users of smaller GPUs
should assess an explicit configuration because lowering \code{nprobe} would
change the selected operating point.

\begin{table}[h]
\centering

\begin{tabularx}{\textwidth}{@{}Xrrrr@{}}
\toprule
Shape represented in evaluation & $k$ & 0.90 & 0.95 & 0.99 \\
\midrule
Large low-dimensional & 15 & (512, 6) & (512, 6) & (512, 12) \\
Large low-dimensional & 30 & (512, 6) & (1024, 16) & (1024, 32) \\
Large low-dimensional & 50 & (1966, 90) & (1966, 90) & (1966, 90) \\
Large low-dimensional & 100 & (5898, 308) & (5898, 308) & (5898, 308) \\
ImageNet features & 15 & (1590, 160) & (1590, 160) & Flat \\
ImageNet features & 30 & (1060, 99) & (1060, 99) & Flat \\
ImageNet features & 50 & (1590, 160) & (1590, 160) & Flat \\
ImageNet features & 100 & (1590, 160) & (1590, 160) & Flat \\
\bottomrule
\end{tabularx}
\caption{Selected CUDA IVF parameters as $(\code{nlist},\code{nprobe})$.
``Flat'' indicates a target-dependent route change.}
\label{tab:auto-ivf-parameters}
\end{table}

\subsection{Regret and runtime comparisons}

Table~\ref{tab:selector-regret-dataset} disaggregates the regret analysis in
the main article by dataset and compares the installed policy with a post hoc
cross-fitted reconstruction.

The two analyses use different candidate sets. The installed policy chose
Flat or IVF, whereas the complete explicit routes available for
cross-fitting were exact-family methods and CAGRA; explicit IVF was absent.
The cross-fitted results are therefore a candidate-set sensitivity analysis.
The table's net time difference is
$T_{\mathrm{exact}}-T_{\mathrm{selected}}$, summed across the 36
metric--$k$--target combinations for each dataset. It describes this benchmark,
not savings weighted by an external application's workload.

\begin{table}[h]
\centering

\small
\setlength{\tabcolsep}{4pt}
\begin{tabular}{@{}lrrr@{\hspace{5mm}}rrr@{}}
\toprule
& \multicolumn{3}{c}{Experimental installed policy} & \multicolumn{3}{c}{Post hoc sensitivity} \\
\cmidrule(lr){2-4}\cmidrule(lr){5-7}
Dataset & \shortstack{Median\\regret} & \shortstack{Approximate\\faster} & \shortstack{Net\\seconds} & \shortstack{Median\\regret} & \shortstack{Approximate\\faster} & \shortstack{Net\\seconds} \\
\midrule
COIL20 & 1.14 & -- & $-0.6$ & 1.00 & -- & $2.4$ \\
FashionMNIST & 1.02 & -- & $-0.7$ & 1.00 & -- & $0.2$ \\
FR-FCM-ZYRM & 1.63 & 12/12 & $3317.7$ & 1.00 & 9/9 & $6561.5$ \\
MNIST & 1.01 & -- & $-0.4$ & 1.00 & -- & $0.3$ \\
MetRef & 1.48 & -- & $-0.7$ & 1.00 & -- & $2.5$ \\
USPS & 1.33 & -- & $-0.5$ & 1.00 & -- & $2.5$ \\
flow18 & 1.88 & 12/12 & $141.2$ & 1.00 & 12/12 & $290.4$ \\
ImageNet features & 1.00 & 0/8 & $-8001.1$ & 1.00 & -- & $33.1$ \\
mass41 & 1.88 & 12/12 & $128.8$ & 1.00 & 12/15 & $264.2$ \\
\bottomrule
\end{tabular}
\caption{CUDA selector regret and absolute time difference by dataset under the
cold full-self-search workload.}
\label{tab:selector-regret-dataset}
\end{table}

Cell-level regret and route-choice records are available in the
\href{https://github.com/tkcaccia/faissR/tree/main/manuscript/jss}{public
GitHub replication materials}. They describe one workload,
\code{cold\_full\_self\_search}. The 324 cells are nested within nine datasets
and are not 324 independent experimental units.

The complete CUDA comparison contained automatic method selection, explicit exact
Flat, brute force, and CAGRA, each with 1,944 successful executions. Every
timed call used \code{nn()} and returned host-resident matrices. The CPU
explicit HNSW set also contained 1,944 successful executions.
Table~\ref{tab:supp-cuda-paired} summarizes CUDA ratios across datasets;
Table~\ref{tab:supp-cuda-dataset} shows their heterogeneity. Among the 44
IVF selections, 36 were faster than explicit exact search. The median
dataset-level ratio for IVF-selected cells was 4.59, but the ImageNet
ratio of 0.14 favored exact search.

\begin{table}[ht]
\centering
\small

\begin{tabularx}{\textwidth}{@{}Xlrrr@{}}
\toprule
Comparison & Datasets & Median & IQR & Range \\
\midrule
Automatic method / explicit exact & 9 & 0.99 & 0.98--1.00 & 0.94--1.01 \\
Automatic method / brute force & 9 & 0.87 & 0.83--1.18 & 0.70--1.90 \\
Automatic method / CAGRA & 9 & 2.23 & 1.00--2.58 & 0.81--10.96 \\
IVF-selected / explicit exact & 4 & 4.59 & 2.80--5.58 & 0.14--5.83 \\
\bottomrule
\end{tabularx}
\caption{Secondary matched CUDA end-to-end ratios. Each ratio is
\(T_{\mathrm{second\ route}}/T_{\mathrm{automatic}}\), so values above one
favor automatic method selection. Dataset medians are summarized across datasets.}
\label{tab:supp-cuda-paired}
\end{table}

\begin{table}[ht]
\centering
\small

\setlength{\tabcolsep}{3.2pt}
\begin{tabular}{@{}lrrrrrr@{}}
\toprule
Dataset & Automatic target & Exact/automatic & Brute/automatic & CAGRA/automatic & IVF target & Exact/IVF \\
\midrule
COIL20 & 36/36 & 0.98 & 0.87 & 1.00 & 0/0 & -- \\
FashionMNIST & 36/36 & 0.98 & 1.18 & 10.94 & 0/0 & -- \\
FR-FCM-ZYRM & 36/36 & 1.00 & 1.00 & 0.81 & 12/12 & 3.69 \\
MNIST & 36/36 & 0.99 & 1.18 & 10.96 & 0/0 & -- \\
MetRef & 36/36 & 0.94 & 0.70 & 2.40 & 0/0 & -- \\
USPS & 36/36 & 0.96 & 0.75 & 2.58 & 0/0 & -- \\
flow18 & 36/36 & 1.00 & 0.83 & 0.96 & 12/12 & 5.83 \\
ImageNet features & 36/36 & 1.01 & 1.90 & 2.23 & 8/8 & 0.14 \\
mass41 & 36/36 & 1.00 & 0.85 & 1.00 & 12/12 & 5.49 \\
\bottomrule
\end{tabular}
\caption{CUDA results by dataset. Automatic-method validation attainment covers 36
metric--$k$--target cells. All routes completed without timeout. Ratios use
second-route time divided by automatic-route time. The last ratio is restricted
to cells in which automatic method selection chose IVF; values above one favor IVF.
Ratio entries are dataset medians of matched replicate ratios.}
\label{tab:supp-cuda-dataset}
\end{table}

\subsection{Dataset and domain holdouts}

The leave-one-dataset-out (LOODO) reconstruction removed every calibration
row belonging to the held-out dataset. It selected a method using the same
prespecified shape group where possible, or the three nearest training
datasets in $\log(n)$--$\log(p)$ space otherwise. The omitted dataset was
used only for evaluation. CUDA selected an exact-audited family in 288 cells
and CAGRA in 36, with no abstentions. Median regret was 1.00, the 90th
percentile was 1.00, and the range was 1.00--2.24. Method-family agreement
with the empirical oracle was 300/324. Pairwise and joint comparisons with
the installed policy are available in the
\href{https://github.com/tkcaccia/faissR/tree/main/manuscript/jss}{public
GitHub replication materials}.

Holding out an entire domain gave less favorable results
(Table~\ref{tab:supp-domain-holdout}). In particular, the cytometry holdout
had median regret 3.08 and a 95th percentile of 6.93. This shows why removing
one named dataset may be a weak test when related datasets remain in the
collection. The grouped analysis covers three domains and the same explicit
candidate set; its scope ends at the observed hardware, geometries, and shape
range.

\begin{table}[ht]
\centering
\small

\begin{tabular}{@{}lrrrrrr@{}}
\toprule
Held-out domain & Datasets & Cells & Target met & Oracle agreement & Median & 95th percentile \\
\midrule
Image-derived & 5 & 180 & 176 & 153 & 1.06 & 6.19 \\
Cytometry & 3 & 108 & 108 & 66 & 3.08 & 6.93 \\
Metabolomics & 1 & 36 & 36 & 36 & 1.00 & 1.00 \\
\bottomrule
\end{tabular}
\caption{Grouped leave-one-domain-out sensitivity reconstructed from the
available explicit independent-query CUDA routes. Regret is selected-route time divided
by empirical-oracle time.}
\label{tab:supp-domain-holdout}
\end{table}

The analogous CPU reconstruction attained 301/324 operating points and
abstained in four. It agreed with the oracle method family in 292/313
evaluable cells. HNSW was selected in 309 cells and exact search in 15;
11 exact selections passed the held-out exactness audit. Median regret was
1.00, the 90th percentile was 1.00, and the range was 1.00--3.16.
These results apply to the reconstructed policy; installed CPU
\code{method = "auto"} was evaluated separately. The replication directory
\path{cpu_loodo/} contains all 324
rows, dataset and metric summaries, the audit report, and checksum ledger.

\section{Replication}

The replication bundle provides the source package, reference manual,
vignette, practical CPU example, and a single commented article script with
HyperText Markup Language (HTML) output. It includes checksummed results,
software and dataset fingerprints, and the files needed to reconstruct the
reported tables.

The script has three modes. \code{compact} runs the interface example on an
ordinary CPU computer. \code{archive} verifies checksums before extraction
and reconstructs numerical summaries from the stored results.
\code{all} runs both. Archive reanalysis runs on CPU; repeating the CUDA
experiments requires compatible NVIDIA hardware and the recorded
FAISS/cuVS build. A checksum, schema, dataset-fingerprint, aggregation, or
table-audit failure stops execution. A successful run records verification
results, per-table SHA-256 digests, and a fresh \code{sessionInfo.txt}.

The archive table builder produces 18 numerical summaries with stable file
identifiers, which are not article table numbers. Separate checksum-verified
inputs reconstruct the controlled CPU, seven-package, tuned HNSW, and
query-workload summaries. Interface and installation tables summarize software
contracts and tested environments.
The bundle is a checksummed snapshot; a persistent archival identifier is
not yet available.

\bibliographystyle{jss}
\bibliography{faissR_jss}

\end{document}
